\documentclass[10pt]{article}
\usepackage[utf8]{inputenc}
\usepackage[T1]{fontenc}
\usepackage{amsmath}
\usepackage{amsfonts}
\usepackage{amssymb}
\usepackage[version=4]{mhchem}
\usepackage{stmaryrd}
\usepackage{graphicx}
\usepackage[export]{adjustbox}
\graphicspath{ {./images/} }

\DeclareUnicodeCharacter{0131}{$\imath$}

\begin{document}
в силу гидродинамич. формулы Стокса), а белый шум $W^{\prime}(t)$ - это «случайная сила», порожденная хаотич. толчками молекул среды, находящихся в тепловом движении, и являющаяся основной причиной броуновского движения. В первоначальной теории броуновского движения, развитой А. Эйнштейном (A. Einstein) и M. Смолуховским (M. Smoluchowski) в 1905-06, пренебрегалось инерцией частıцы, т. е. считалось, что $m=0$; при этом уравнение (*) приводило к выводу, что координата броуновской частицы

\[
X(t)=\int_{0}^{t} V\left(t^{\prime}\right) d t^{\prime}
\]

равна $\beta^{-1} W(t)$, т. е. представляет собой вıнеровский процесс. Таким образом, винеровский процесс описывает модель Эйнштейна - Смолуховского броуновского движения (отсюда другое его название - п р о ц е с с бр о унов с ко г о д в и же ни я); т. к. этот процесс недифференцируем, то в теории Эйнштейна - Смолуховского частица, совершающая броуновское движение, не имеет конечной скорости. Уточненная теория броуновского движения, опирающаяся на уравнение $(*)$, где $m \neq 0$, была предложена Л. Орнштейном и Дж. Уленбеком ([1]; см. также [2]); позже та же теория была выдвинута С. Н. Бернштейном [3] и А. Н. Колмогоровым [4]. В теории Орнштейна - Уленбека скорость $V(t)$ броуновской частицы является конечной, но ее ускорение бесконечно (так как О.-У. п. недиффференцируем); для того чтобы и ускорение оказалось конечным, надо уточнить теорию, учтя отличие случайной силы от идеализированного белтого шума $W^{\prime}(t)$.

Уравнение (*) можно использовать и для описания одномерного броуновского движения гармонич. осциллятора, если пренебречь его массой и считать, что $V(t)$ - это координата осциллятора, - $\frac{m d V}{d t}-$ сила вязкого трения, $-\beta V-$ регулярная упругая сила, удерживающая осциллятор, а $W^{\prime}(t)$ - случайная сила, создаваемая молекулярными толчками. Таким образом, О.-У. п. доставляет также модель пульсаций координаты гармонич. осциллятора, совершающего броуновское движение, родственную модели Эйнштейва - Смолуховского броуновского движения свободної частицы.

О. - У. П. является однородным по времени марковским процессом дифффузионного типа (см. Диффуззионный процесc); наоборот, процесс $V(t)$, являющийся одновременно стационарным случайным процессом, гауссовским процессом и марковским процессом, обязательно представляет собой О.- У. П. Как марковский процесс О.- У. п. удобно характеризовать его переходної плотностью вероятности $p(t, x, y)$, представляющей собой фундаментальное решение соответствуюшего уравнения Фоккера - Планка (т. е. прямого Колмогорова уравнения) вида

\[
\frac{\partial p}{\partial t}=\alpha \frac{\partial(y p)}{\partial y}+\alpha \sigma^{2} \frac{\partial^{2} p}{\partial y^{2}},
\]

и, следовательно, задаваемой формулой

\[
p(t, x, y)=\frac{1}{\left[2 \pi \sigma^{2}\left(1-e^{-2 \alpha t}\right)\right]^{2}} \exp \left\{-\frac{\left(y-x e^{-\alpha t}\right)^{2}}{2 \sigma^{2}\left(1-e^{-2 \alpha t}\right)}\right\} .
\]

Многие свойства О.- У. п. $V(t)$ (включая и его марковость) можно вывести іІз известных свойств винеровского процесса, воспользовавшись тем, что процесс

\[
W_{0}(t)=\frac{\sqrt{t}}{\sigma} V\left(\frac{\ln t}{2 \alpha}\right)
\]

является стандартным винеровским процессом (см. [5]) В частности, отсюда следует, что реализации О.- У. II. непрерывны і нигде не дифференцируемы с вероятностью 1 и что

\[
\varlimsup_{t \rightarrow 0} \frac{|V(t)-V(0)|}{\sqrt{4 \alpha \sigma^{2} t \ln \ln \frac{1}{t}}}=1, \varlimsup_{t \rightarrow \infty} \frac{|V(t)|}{\sqrt{2 \sigma^{2} \ln t}}=1
\]

с вероятностью 1.

Jum.: [1] U h l e n h e c k G.E., O rnste in L. S., "Phys. Rev.», 1930, v. 36, p. 823-41; [2] ч a н д р a c e к a p C., Сто-
хастические пгроблемы в физике и астрономии, пер. с англ., М., 1947; [3] Б е р н ш т е й н С. Н., "Докл. АН СССР", 1934, т. 1 , № 1 , c. $1-9 ; \mathcal{N}$ 7 7 c. $361-65 ;[4] \mathrm{K}$ o $1 \mathrm{~m}$ o g o r o v A. N., "Ann. Math.», 1934, v. 35, p. 116-17; [5] D o o b J. L., "AAnn.
Math.», 1942, v. 43, p. 35i-69. A2лом. ОРРА - ЗОММЕРФЕЛЬДА УРАВНЕНИЕ - линеЙное обыкновенное дифференциальное уравнение

$\varphi^{(4)}-2 \alpha^{2} \varphi^{\prime \prime}+\alpha^{4} \varphi=i \alpha R\left[(w-c)\left(\varphi^{\prime \prime}-\alpha^{2} \varphi\right)-w^{\prime \prime} \varphi\right], \quad$ (1) где $R$ - число Рейнольдса, $w(y)$ - заданная фуунция (профиль скорости невозмущенного потока), к-рая обычно предполагается голоморфной в окрестность отрезка $[-1,1]$ в комплексной плоскости $y, \alpha>0-$ постоянная и $c$ - спектральный параметр. Для О.3. у. исследуется краевая задача

\[
\varphi(-1)=\varphi^{\prime}(-1)=\varphi(1)=\varphi^{\prime}(1)=0 \text {. }
\]

О. - 3. у. возникло при исследовании У. Орром [1] и А. Зоммерфельдом [2] устойчивости в линейном прпближении плоского течения Пуазёйля - течения вязкой несжимаемой жидкости в слое $-\infty<x<\infty,-1<$ $<y<1$ с твердыми границами; возмущение для функции тока берется в виде $\varphi(y) e^{i \alpha(x-c t)}$.

Собственные значения задачи (1), (2), вообще говоря, комплексны; течение устойчиво, если Im $c<0$ для всех собственных значений, и неустойчиво, если $\operatorname{Im} c>0$ для нек-рого из них. Кривая $\operatorname{Im} c(\alpha, R)=0$ наз. н е й т р а Л ьн о й к р и в о й. Течение Пуазёйля устойчиво при небольших числах Рейнольдса. В. Гейзенберг [6] впервые высказал предположение, что течение Пуазёйля неустойчиво при болыших числах Рейнольдса, и вычислил 4 точки нейтральной кривой. Для квадратичного профиля скорости установлено, что течение неустойчиво при $\alpha R \gg 1$.

Асимптотич. теория 0.- 3. у. построена в предположении, что $(\alpha R)^{-1} \rightarrow 0-$ малый параметр. Точка $y_{c}$, в к-рой $w\left(y_{c}\right)=0$, является точкой поворота (см. Малого параметра метод). В малой окрестности точки $y \neq y_{c}$ О.- 3. у. имеет фундаментальную систему решөний вида

\[
\begin{gathered}
\varphi_{1,2}(y)=\varphi_{1,2}^{0}(y)+O\left((\alpha R)^{-1}\right), \\
\varphi_{3,4}(y)=\exp \left[ \pm \int^{y} \sqrt{i(w-c)} d y\right] \times \\
\times\left[(w-c)^{-5 / 4}+O\left((\alpha R)^{-1 / 2}\right)\right],
\end{gathered}
\]

где $\varphi_{1}^{0}(y), \varphi_{2}^{0}(y)$ - фундаментальная система решений невязкого (то есть $\alpha R=0$ ) уравнения

\[
(w-c)\left(\varphi^{\prime \prime}-\alpha^{2} \varphi\right)-w^{\prime \prime} \varphi=0 .
\]

Исследование задач (1), (2) связано, вапр., со следующими труцностями: 1) невязкое уравнение в окрестности точки $y=y_{c}$ имеет голоморфное в ней решение и решөние с логарифмич. особенностью; 2) при малых $|c|$ (т. е. в наиболее важном случае) точки поворота сливаются с концами отрезка $[-1,1]$ (напр., для квадратичного профиля скорости $w=1-y^{2}$ ).

При $\alpha R \gg 1$ получено строгое обоснование неустойчивости (см. [3], [4]).

Jum.: [1] O r r W. M c F., "Proc. R. Irish. Acad. A», 1907. del IV Congresso internazionale del matematici (Roma, 1908), 1909, p. $116-24$; [3] Л и н ь Ц з я- ऍ 3 я о, Теория гидродинамической устойчивости, пер. с англ., М., 1958; 14] Гидродинамическая неустойчивость. Cборник, пер. с англ., M., 1964; [5] G e r-
s t ing J. M., J a n o w s i D. F., "Internat. J. Num. Meth. in Eng.", 1972, v. 4, p. 195-206; [6] H e i s e n ber g W., "Ann.

\begin{center}
\includegraphics[max width=\textwidth]{2023_07_08_7ecd5aa9d848744d53aag-1}
\end{center}


\end{document}